\section{Objetivo general}
Diseñar y fabricar una impresora 3D basada en tecnología FDM orientada a aplicaciones académicas, mediante el rediseño del sistema de extrusión y enfriamiento de capa, con el propósito de mejorar la calidad de impresión y evaluar su desempeño frente a equipos comerciales equivalentes.

\section{Objetivos específicos}
\begin{itemize}
  \item Diseñar la estructura mecánica y los componentes del sistema de impresión 3D mediante software CAD.
  \item Fabricar las piezas estructurales y funcionales utilizando procesos de manufactura aditiva.
  \item Integrar los sistemas mecánicos, electrónicos y térmicos de la impresora 3D.
  \item Configurar y calibrar el firmware de control para garantizar el correcto funcionamiento del equipo.
  \item Evaluar el comportamiento térmico del hotend y del sistema de enfriamiento de capa durante el proceso de impresión.
  \item Analizar la precisión dimensional y calidad superficial de las piezas fabricadas.
  \item Comparar el desempeño técnico y el costo del sistema desarrollado con impresoras comerciales de características similares.
  \item Validar la aplicabilidad del equipo como herramienta didáctica para entornos educativos y de prototipado.
\end{itemize}

\section{Metas}
Para dar cumplimiento a los objetivos planteados, se establecieron las 
siguientes metas concretas y medibles

\begin{enumerate}

\item Obtener una impresora 3D funcional basada en arquitectura abierta y tecnología FDM, lista para la impresión de piezas de prueba.

\item Lograr una reducción significativa en la masa del módulo de extrusión 
para disminuir la inercia sobre el carro del eje \(X\) y mejorar la dinámica 
del sistema.

\item Mejorar la disipación térmica del hotend y garantizar la integridad de 
los componentes del extrusor, evitando fallas por sobrecalentamiento o 
deformación.

\item Reducir la aparición de defectos de impresión como \textit{stringing}, \textit{warping} y \textit{ghosting} mediante la optimización de parámetros mecánicos y térmicos.

\item Obtener piezas con precisión dimensional y calidad superficial adecuadas para aplicaciones didácticas y de prototipado funcional.

\item Implementar un equipo con arquitectura abierta que permita futuras 
mejoras mecánicas, electrónicas y de firmware, así como la fácil sustitución 
de componentes.

\item Demostrar la viabilidad técnica y económica de la impresora, logrando 
un costo de fabricación inferior al de las alternativas comerciales 
equivalentes de gama baja.


\end{enumerate}


